\subsection{Summary of the ATC Decoupling Framework}

This chapter has established the distributed coordination layer of the Multi-Microgrid system through the Analytical Target Cascading (ATC) framework. To overcome the computational bottlenecks and severe privacy vulnerabilities inherent in centralized dispatch architectures, ATC implements a strict spatial decoupling strategy. Unlike flat consensus mechanisms such as ADMM, ATC perfectly aligns with the natural hierarchical topology of distribution grids, enabling secure interaction between the central distribution operator (DSO) and autonomous Microgrids via a boundary Target-Response mechanism.

Mathematically, global consensus is enforced without exposing sensitive internal constraints by embedding an Augmented Lagrangian penalty function into the local SOCP objectives. By employing linear shadow prices ($\lambda$) and quadratic penalty multipliers ($\rho$) coupled with adaptive residual balancing, the algorithm guarantees rapid and stable convergence. While this chapter successfully formulates the spatial coordination of the network, grid operations remain inherently susceptible to real-time physical disturbances and stochastic uncertainties. Consequently, this static hierarchical foundation serves as the critical prerequisite for integrating temporal dynamism. In the subsequent chapter, this ATC framework will be augmented with a Model Predictive Control (MPC) rolling-horizon mechanism to establish a fully dynamic, proactive resilience architecture capable of real-time graceful degradation.
